\documentclass[11pt]{article}

\usepackage[margin=0.85in]{geometry}
\usepackage{amsmath,amssymb}
\usepackage{enumitem}
\usepackage{xcolor}
\usepackage{tcolorbox}
\usepackage{fancyhdr}
\usepackage{lastpage}
\usepackage{hyperref}
\usepackage{microtype}

\definecolor{uwpurple}{HTML}{4B2E83}
\definecolor{uwgold}{HTML}{B7A57A}
\definecolor{softpurple}{HTML}{F4F1FA}

\hypersetup{
    colorlinks=true,
    linkcolor=uwpurple,
    urlcolor=uwpurple
}

\pagestyle{fancy}
\fancyhf{}
\lhead{\textbf{MATH 394: Probability I}}
\rhead{\textbf{Homework 4}}
\cfoot{\thepage\ of \pageref{LastPage}}

\setlength{\parindent}{0pt}
\setlength{\parskip}{0.45em}

\newcommand{\course}{MATH 394: Probability I}
\newcommand{\instructor}{Arman Jahangiri}
\newcommand{\department}{Department of Mathematics}
\newcommand{\university}{University of Washington}

\newtcolorbox{problemheader}{
    colback=softpurple,
    colframe=uwpurple,
    boxrule=0.8pt,
    arc=2mm,
    left=2mm,
    right=2mm,
    top=1mm,
    bottom=1mm
}

\newcommand{\source}[1]{%
    \vspace{0.15cm}
    {\small\textit{Source: Blitzstein and Hwang, \textit{Introduction to Probability}, Chapter 3, Problem #1.}}
}

\newcommand{\answerbox}{%
\begin{flushright}
\begin{tcolorbox}[
    width=0.36\textwidth,
    height=2cm,
    colback=white,
    colframe=uwpurple,
    boxrule=0.7pt,
    arc=1mm
]
{\small\textbf{Final Answer}}
\end{tcolorbox}
\end{flushright}
}

\newcommand{\partanswerbox}[1]{%
\begin{flushright}
\begin{tcolorbox}[
    width=0.36\textwidth,
    height=2cm,
    colback=white,
    colframe=uwpurple,
    boxrule=0.7pt,
    arc=1mm
]
{\small\textbf{Final Answer for Part #1}}
\end{tcolorbox}
\end{flushright}
}

\begin{document}

\begin{titlepage}
\centering
\vspace*{1.2cm}

{\Large \textbf{\university}}\\
{\large \department}

\vspace{1.5cm}

{\Huge \color{uwpurple}\textbf{Homework 4}}\\[0.25cm]
{\Large \textbf{\course}}\\[0.25cm]
{\large Instructor: Arman Jahangiri}

\vspace{1cm}

\begin{tcolorbox}[
    width=0.78\textwidth,
    colback=softpurple,
    colframe=uwpurple,
    boxrule=0.9pt,
    arc=3mm
]
\centering
\textbf{Submission:} Single PDF on Gradescope

\textbf{Deadline:} 10:00 PM Pacific Time on the listed due date
\end{tcolorbox}

\vspace{0.8cm}

\begin{tcolorbox}[
    width=0.86\textwidth,
    colback=white,
    colframe=uwgold,
    boxrule=0.8pt,
    arc=3mm
]
This homework is worth \textbf{50 points}.

Across the quarter, there are \textbf{eight homework assignments}, worth \textbf{400 points total}. Homework assignments together account for \textbf{40\%} of the final course grade. Further course policies can be seen in the \href{https://armanjg.github.io/MATH-394-Probability-1/}{\textbf{MATH 394 syllabus}.}
\end{tcolorbox}

\vfill
\end{titlepage}

\section*{Homework Policy}

\subsection*{Submission and Deadline}

Please:

\begin{itemize}[leftmargin=*]
    \item Submit homework as a single PDF on Gradescope by \textbf{10:00 PM (Pacific Time)} on the listed due date.
    \item Begin each problem on a new page, and ensure that pages are correctly tagged in Gradescope.
    \item Write solutions clearly and legibly.
    \item Typesetting solutions in LaTeX is encouraged, but not required.
\end{itemize}

\subsection*{Late Days}

Each student is allotted \textbf{six late days} for the quarter. A late day extends a homework deadline by up to 24 hours without penalty. For example, submitting an assignment anytime between 10:01 PM on the due date and 10:00 PM the following day counts as one late day.

The following rules apply:
\begin{itemize}[leftmargin=*]
   \item No explanation is required to use late days.
   \item At most \textbf{two late days} may be used on any single homework assignment.
   \item Late days may not be used on the final homework assignment or on any assignment for which solutions have already been released.
   \item It is the student's responsibility to keep track of how many late days have been used during the quarter.
\end{itemize}

Once all late days have been exhausted, additional late submissions will incur a penalty of \textbf{10\% per day}, up to a maximum deduction of \textbf{50\%}. Assignments submitted more than five days late, or after solutions have been released, will not be accepted.

\subsection*{Submission Issues and Technical Difficulties}

If a serious technical issue prevents a timely Gradescope submission, students may temporarily submit their assignment by email to the instructor at \href{mailto:armanjg@uw.edu}{armanjg@uw.edu}. In such cases:

\begin{itemize}[leftmargin=*]
    \item The submission must consist of a \textbf{single PDF file}. The email subject line must follow the format:
    \[
    \texttt{HW\# Submission - FULL NAME - STUDENT ID}.
    \]
    \item Once the Gradescope issue has been resolved, the student must upload the \emph{exact same file} to Gradescope as soon as possible.
    \item The student should clearly indicate in the Gradescope submission that the assignment was previously submitted by email before the deadline.
\end{itemize}

Email submissions are intended only for genuine technical emergencies and should not be used as a substitute for timely Gradescope submission.

\newpage

% ------------------------------------------------------------

\begin{problemheader}
\textbf{Problem 1. }
\end{problemheader}
\source{1}

People are arriving at a party one at a time. While waiting for more people to arrive, they entertain themselves by comparing their birthdays.

Let $X$ be the number of people needed to obtain a birthday match; that is, before person $X$ arrives, no two people have the same birthday, but when person $X$ arrives there is a match.

Find the PMF of $X$.

\vfill

\newpage

% ------------------------------------------------------------

\begin{problemheader}
\textbf{Problem 2. }
\end{problemheader}
\source{3}

Let $X$ be a random variable with CDF $F$, and let
\[
Y=\mu+\sigma X,
\]
where $\mu$ and $\sigma$ are real numbers with $\sigma>0$.

The random variable $Y$ is called a \emph{location--scale transformation} of $X$.

Find the CDF of $Y$ in terms of $F$.

\vfill

\newpage

% ------------------------------------------------------------

\begin{problemheader}
\textbf{Problem 3. }
\end{problemheader}
\source{4}

Let $n$ be a positive integer and define
\[
F(x)=\frac{\lfloor x\rfloor}{n}
\]
for $0\le x\le n$, where $\lfloor x\rfloor$ is the greatest integer less than or equal to $x$.

Also define
\[
F(x)=0 \quad \text{for } x<0,
\]
and
\[
F(x)=1 \quad \text{for } x>n.
\]

\begin{enumerate}[label=(\alph*), leftmargin=*]
    \item Show that $F$ is a valid CDF.

    \item Find the PMF corresponding to $F$.
\end{enumerate}

\vfill


\newpage

% ------------------------------------------------------------

\begin{problemheader}
\textbf{Problem 4. }
\end{problemheader}
\source{5}

\begin{enumerate}[label=(\alph*), leftmargin=*]
    \item Show that
    \[
    p(n)=\left(\frac12\right)^{n+1},
    \qquad n=0,1,2,\dots,
    \]
    is a valid PMF for a discrete random variable.

    \item Find the CDF of a random variable with the PMF from part (a).
\end{enumerate}

\vfill


\newpage

% ------------------------------------------------------------

\begin{problemheader}
\textbf{Problem 5. }
\end{problemheader}
\source{8}

There are 100 prizes, with one worth \$1, one worth \$2, \dots, and one worth \$100. There are 100 boxes, each containing one of the prizes.

You get 5 prizes by picking random boxes one at a time, without replacement.

Find the PMF of how much your most valuable prize is worth, as a simple expression in terms of binomial coefficients.

\vfill


\newpage

% ------------------------------------------------------------

\begin{problemheader}
\textbf{Problem 6. }
\end{problemheader}
\source{9}

Let $F_1$ and $F_2$ be CDFs, let $0<p<1$, and define
\[
F(x)=pF_1(x)+(1-p)F_2(x)
\]
for all $x$.

\begin{enumerate}[label=(\alph*), leftmargin=*]
    \item Show directly that $F$ has the properties of a valid CDF. The distribution defined by $F$ is called a \emph{mixture} of the distributions defined by $F_1$ and $F_2$.

    \item Consider creating a random variable in the following way. Flip a coin with probability $p$ of Heads. If the coin lands Heads, generate a random variable according to $F_1$; if the coin lands Tails, generate a random variable according to $F_2$. Show that the random variable obtained in this way has CDF $F$.
\end{enumerate}

\vfill


\newpage

% ------------------------------------------------------------

\begin{problemheader}
\textbf{Problem 7.}
\end{problemheader}
\source{22}

There are two coins, one with probability $p_1$ of Heads and the other with probability $p_2$ of Heads. One of the coins is randomly chosen, with equal probabilities for the two coins. It is then flipped $n\ge2$ times.

Let $X$ be the number of times it lands Heads.

\begin{enumerate}[label=(\alph*), leftmargin=*]
    \item Find the PMF of $X$.

    \item What is the distribution of $X$ if $p_1=p_2$?

    \item Give an intuitive explanation of why $X$ is not Binomial for $p_1\ne p_2$. You may assume that $n$ is large for your explanation.
\end{enumerate}

\vfill



\newpage

% ------------------------------------------------------------

\begin{problemheader}
\textbf{Problem 8. }
\end{problemheader}
\source{30}

A certain company has $n+m$ employees, consisting of $n$ women and $m$ men. The company is deciding which employees to promote.

\begin{enumerate}[label=(\alph*), leftmargin=*]
    \item Suppose for this part that the company decides to promote $t$ employees, where
    \[
    1\le t\le n+m,
    \]
    by choosing $t$ random employees, with equal probabilities for each set of $t$ employees.

    What is the distribution of the number of women who get promoted?

    \item Now suppose that instead of having a predetermined number of promotions to give, the company decides independently for each employee, promoting the employee with probability $p$.

    Find the distributions of:
    \begin{itemize}
        \item the number of women who are promoted,
        \item the number of women who are not promoted,
        \item the number of employees who are promoted.
    \end{itemize}

    \item In the set-up from part (b), find the conditional distribution of the number of women who are promoted, given that exactly $t$ employees are promoted.
\end{enumerate}

\vfill



\end{document}